% fig1_overview — 方法总体流程（S1–S4）
% 用途: 交底书 图1, §五 技术方案总体流程 (对应权利要求方法步骤顺序)
% 规范: 全中文标注; 黑白线条; 图内只保留必需词语(标题/说明见 FIGURES.md 图注)
% 编译: bash docs/patent/build.sh   (xelatex + ctex)
\documentclass[tikz,border=8pt]{standalone}
\usepackage[fontset=windows]{ctex}
\usepackage{amsmath}
\usepackage{tikz}
\usetikzlibrary{arrows.meta,positioning}
\begin{document}
\begin{tikzpicture}[
  font=\small,
  io/.style={draw,rounded corners=2pt,align=center,text width=11.2cm,fill=gray!10},
  box/.style={draw,rounded corners=2pt,align=center,text width=11.2cm,fill=white},
  result/.style={box,draw=black!75,line width=1.1pt},
  ar/.style={-{Stealth[length=3mm]},thick},
]
\node[io] (inp) {输入：目标点集合（坐标、需求量）、仓库位置、无人机参数（$\alpha$、$\beta$、载重上限、电池容量）};

\node[box,below=0.72cm of inp] (s1) {S1\quad 载重感知等效电量距离（逐段按出发载重求值）\\[2pt]
  $\mathrm{EED}(i,j,L)=d_{ij}\cdot\dfrac{\alpha+\beta L}{\alpha}$\quad（$L$ 为出发时载重）};

\node[box,below=0.72cm of s1] (s2) {S2\quad 电量感知最近邻构造（N 起点）\\[2pt]
  载重预检：总需求 $>$ 载重上限则判定不可行；每步要求往返电量可行};

\node[box,below=0.72cm of s2] (s3) {S3\quad 增量局部搜索：2-opt 与 Or-opt 交替\\[2pt]
  首次改进接受；每次候选移动仅按受影响段做增量电量校验（载重由预检保证）};

\node[box,below=0.72cm of s3] (s4) {S4\quad 逐段全量后验证（载重与能耗）};

\node[result,below=0.72cm of s4] (res) {输出：访问顺序、总等效电量距离、可行性};

\draw[ar] (inp)--(s1); \draw[ar] (s1)--(s2); \draw[ar] (s2)--(s3); \draw[ar] (s3)--(s4); \draw[ar] (s4)--(res);
\end{tikzpicture}
\end{document}
